\subsection{An actual unitary automorphic realization of one cubic plethysm ring}

There is a precise automorphic dictionary in a fixed dimension that is
stronger than an analogy between two sign arguments. We use the following
known theorem with its exact grading, and then transfer the pole already
proved here.

\paragraph{Known dictionary.}
Looijenga--Swierstra, \emph{The period map for cubic threefolds},
arXiv:\texttt{math/0608279v1}, Theorem~3.1, printed p.~12, identifies
$\mathbb C[\operatorname{Sym}^3(\mathbb C^5)^*]^{SL_5}$ with the
ring of meromorphic automorphic forms on their arithmetic complex
$10$-ball quotient, allowing poles along the specified hyperplane
arrangement. Write $\mathcal M_k$ for its weight-$k$ space.
The isomorphism multiplies automorphic weight by $5/3$.
The central-character explanation on p.~11 gives
$\mathcal M_k=0$ unless $3\mid k$; negative weights vanish
by the theorem. The Hermitian lattice and automorphic line
bundle are defined on pp.~10--11. This is a known ring
isomorphism, including multiplication, rather than a
new construction of an automorphic model.

\begin{corollary}[P-F37: exact automorphic regrading and inherited poles]
\label{cor:automorphiccubic}
With those meromorphic spaces and their multiplication,
\[
 \bigoplus_{s\ge0}\mathcal M_{3s}\simeq T^{(5)},\qquad
 \dim\mathcal M_{3s}=f_5(s).
 \tag{P-F37a}
\]
The left hand side is graded by $s$ in this formula.
In the original automorphic weight, put
\[
 G(u)=\sum_{k\ge0}\dim\mathcal M_k\,u^k.
\]
Then
\[
                         G(u)=H_5(u^3),
 \tag{P-F37b}
\]
and $G$ has an actual simple pole at every primitive
$48$th root $\eta$. More precisely,
\[
 \lim_{u\to\eta}(1-u^{48})G(u)
                          =\mathcal C_5(\eta^3)\ne0.
 \tag{P-F37c}
\]
Consequently every expression for $G$ with polynomial numerator
and denominator factors $1-u^a$ has some $a$ divisible by $48$.
Every homogeneous system of parameters in this automorphic
ring has some weight at least $48$.
\end{corollary}

\begin{proof}
The cited graded isomorphism sends $\mathcal M_{3s}$
to the ordinary coefficient-degree-$5s$ part of the cubic
invariant ring. P-F19 identifies the latter, with its actual
multiplication and degree $s$, with $T^{(5)}_s$.
Composing the two graded ring isomorphisms proves (P-F37a).
The central-character vanishing outside weights divisible
by three then gives (P-F37b) coefficient by coefficient.

For $n=5$, P-F31 has $M=2^4=16$. If $\eta$ is primitive
of order $48$, then $\eta^3$ is primitive of order $16$.
Substitution $t=u^3$ in P-F31b gives (P-F37c) directly,
since $1-t^{16}=1-u^{48}$. The map $u\mapsto u^3$
has nonzero derivative at $\eta$, so a simple pole
remains simple. Equivalently, the nonzero finite limit
in (P-F37c) and the simple zero of $1-u^{48}$ give
this assertion without any coordinate convention.
If none of the denominator weights were divisible by
$48$, every denominator factor would be nonzero at
$\eta$, contradicting this pole. This proves the
assertion about arbitrary polynomial numerators.

Finally an h.s.o.p. in the automorphic ring, with
its original weights divided by three, is an h.s.o.p.
in $T^{(5)}$ under (P-F37a). All its nonzero weights
are divisible by three. P-F19 forces one of the
regraded parameter degrees to be at least $16$,
so its original automorphic weight is at least $48$.
This proves the last claim independently of choosing
a rational Hilbert presentation.
\end{proof}

\paragraph{What this supplies for the theta-lift question.}
This gives an identified unitary automorphic object carrying
one actual ordinary plethysm algebra and all its stretching
coefficients. It also shows that an automorphic realization
preserves, rather than removes, the grading obstructions proved
in the algebraic model. Its dimension is fixed here; no
uniform family in growing $n$ or index estimate is obtained.

A transfer from the supplied note's theta lifts must still
identify their output with these particular weight spaces,
including the permitted arrangement poles, the change of
weight and multiplication. The ball has unitary signature
$(10,1)$, whereas the supplied note's iterated lifts use
its specified totally definite unitary data. No theorem
identifying these constructions is stated in either the
cited period-map theorem or the supplied \emph{Remark}.
This observation does not preclude an additional comparison,
and does not assert that theta lifts can never have the
needed analytic type. It makes the remaining comparison
concrete in a case where an automorphic dictionary is
already available.

The primary source is
\url{https://arxiv.org/pdf/math/0608279}, Theorem~3.1,
p.~12, with the grading discussion on pp.~10--11.
